\clearpage
\hypertarget{native-481-january10}{}
\begin{center}
\rule{0.55\linewidth}{0.4pt}\\[2pt]
Astronomische Nachrichten. Vol. XXVII. No. 625. pp. 13..16. 29 January 1848.\\[2pt]
\rule{0.55\linewidth}{0.4pt}
\end{center}
\begin{center}[\textsc{Gauss} to \textsc{Schumacher}.]\quad Göttingen, 10 January 1848.\end{center}

I observed Flora on the meridian on the 5th.

\begin{center}
\setlength{\tabcolsep}{4pt}
\begin{tabular}{l|r|r|r}
1848 January 5 & $9^{\rm u}\ 11^{\rm m}\ 44^{\rm s}8$ & $62^\circ\ 30'\ 18''1$ & $+16^\circ\ 9'\ 45''2$ \\
\end{tabular}
\end{center}

The discrepancy in \textsc{Goldschmidt}’s elements is $+16''4$, $+3''5$. On the 3rd I found Flora considerably fainter than in December of last year, but still bright enough that it would probably have stood almost full illumination. On the 5th I had to dim the illumination considerably. In fact, the planet could scarcely stand even low field illumination. I attributed this to the state of the air, since the fixed stars also gave very indistinct and unsteady images. In any case, the observation on the 5th is less accurate than the one on the 3rd.
